\documentclass[tikz,border=14pt]{standalone}
\usepackage[no-math]{fontspec}
\usepackage{xeCJK}
\usepackage[UTF8,fontset=none]{ctex}
\usepackage{tikz}
\usetikzlibrary{arrows.meta,positioning,fit,calc,backgrounds}

\IfFontExistsTF{Times New Roman}{\setmainfont{Times New Roman}}{\setmainfont{TeX Gyre Termes}}
\IfFontExistsTF{Kaiti TC}{\setCJKmainfont[AutoFakeBold=3]{Kaiti TC}}
  {\IfFontExistsTF{Songti TC}{\setCJKmainfont[AutoFakeBold=3]{Songti TC}}
    {\setCJKmainfont{PingFang TC}}}

% thesis_colors.tex ─ 論文統一 TikZ 色票定義
% 在每個 standalone 圖的 preamble 以 % thesis_colors.tex ─ 論文統一 TikZ 色票定義
% 在每個 standalone 圖的 preamble 以 % thesis_colors.tex ─ 論文統一 TikZ 色票定義
% 在每個 standalone 圖的 preamble 以 \input{thesis_colors.tex} 引入

\definecolor{thesisBlue}{RGB}{43,108,176}      % 訓練集、主箭頭、主色
\definecolor{thesisAmber}{RGB}{201,140,30}     % 驗證集、次要強調
\definecolor{thesisCrimson}{RGB}{155,35,53}    % 測試集、警示
\definecolor{thesisTeal}{RGB}{74,157,143}      % 節點類型 1（電子零組件）
\definecolor{thesisForest}{RGB}{45,122,79}     % 節點類型 2（光電及I/O）、LSTM Update
\definecolor{thesisSand}{RGB}{194,168,120}     % 節點類型 3（半導體）
\definecolor{thesisDark}{RGB}{45,55,72}        % 框線、文字、主箭頭
\definecolor{thesisLight}{RGB}{237,242,247}    % 淡背景、節點填色
\definecolor{thesisPurple}{RGB}{107,70,193}    % GAT 第二注意力頭
\definecolor{thesisMidGray}{RGB}{130,145,160}  % 次要線段、虛線
 引入

\definecolor{thesisBlue}{RGB}{43,108,176}      % 訓練集、主箭頭、主色
\definecolor{thesisAmber}{RGB}{201,140,30}     % 驗證集、次要強調
\definecolor{thesisCrimson}{RGB}{155,35,53}    % 測試集、警示
\definecolor{thesisTeal}{RGB}{74,157,143}      % 節點類型 1（電子零組件）
\definecolor{thesisForest}{RGB}{45,122,79}     % 節點類型 2（光電及I/O）、LSTM Update
\definecolor{thesisSand}{RGB}{194,168,120}     % 節點類型 3（半導體）
\definecolor{thesisDark}{RGB}{45,55,72}        % 框線、文字、主箭頭
\definecolor{thesisLight}{RGB}{237,242,247}    % 淡背景、節點填色
\definecolor{thesisPurple}{RGB}{107,70,193}    % GAT 第二注意力頭
\definecolor{thesisMidGray}{RGB}{130,145,160}  % 次要線段、虛線
 引入

\definecolor{thesisBlue}{RGB}{43,108,176}      % 訓練集、主箭頭、主色
\definecolor{thesisAmber}{RGB}{201,140,30}     % 驗證集、次要強調
\definecolor{thesisCrimson}{RGB}{155,35,53}    % 測試集、警示
\definecolor{thesisTeal}{RGB}{74,157,143}      % 節點類型 1（電子零組件）
\definecolor{thesisForest}{RGB}{45,122,79}     % 節點類型 2（光電及I/O）、LSTM Update
\definecolor{thesisSand}{RGB}{194,168,120}     % 節點類型 3（半導體）
\definecolor{thesisDark}{RGB}{45,55,72}        % 框線、文字、主箭頭
\definecolor{thesisLight}{RGB}{237,242,247}    % 淡背景、節點填色
\definecolor{thesisPurple}{RGB}{107,70,193}    % GAT 第二注意力頭
\definecolor{thesisMidGray}{RGB}{130,145,160}  % 次要線段、虛線


\begin{document}
\begin{tikzpicture}[
  arr/.style={-{Stealth[length=1.8mm,width=1.5mm]}, thesisDark!45, line width=0.6pt},
  clustertitle/.style={font=\small\bfseries, thesisDark, anchor=south},
  nodelabel/.style={font=\fontsize{5.8pt}{7pt}\selectfont, thesisDark, align=center},
  centerlabel/.style={font=\footnotesize\itshape, thesisDark!60, align=center},
]


%%── Cluster 1: 電子零組件 ─────────────────────────────────────
\def\cx{0}\def\cy{0}
\def\ndx{1.05}\def\ndy{1.60}
\foreach \i/\ii in {1/tl,2/tr,3/ml,4/mr,5/bl,6/br}{
  \ifnum\i=1 \coordinate (na_\i) at ({-\ndx},{\ndy}); \fi
  \ifnum\i=2 \coordinate (na_\i) at ({+\ndx},{\ndy}); \fi
  \ifnum\i=3 \coordinate (na_\i) at ({-\ndx},{0});    \fi
  \ifnum\i=4 \coordinate (na_\i) at ({+\ndx},{0});    \fi
  \ifnum\i=5 \coordinate (na_\i) at ({-\ndx},{-\ndy});\fi
  \ifnum\i=6 \coordinate (na_\i) at ({+\ndx},{-\ndy});\fi
}
% Edges
\foreach \i in {1,...,6}{
  \foreach \j in {1,...,6}{
    \ifnum\i<\j
      \draw[arr] (na_\i) -- (na_\j);
    \fi
  }
}
% Nodes
\foreach \k/\lbl in {1/6133\\金橋,2/6224\\英典,3/8210\\勤誠,4/6282\\康舒,5/3015\\德淵電,6/6235\\華宇}{
  \node[circle, draw=thesisTeal!80!thesisDark, fill=thesisTeal!30, thick, minimum size=10mm]
       (Na_\k) at (na_\k) {};
  \draw[arr, looseness=7, out=105, in=75] (Na_\k) to (Na_\k);
  \node[nodelabel, below=5pt] at (Na_\k.south) {\lbl};
}
\node[clustertitle] at (0, {1.35+0.75}) {電子零組件};

%%── Cluster 2: 光電及I/O ──────────────────────────────────────
\begin{scope}[xshift=5.2cm]
\foreach \i in {1,...,6}{
  \ifnum\i=1 \coordinate (nb_\i) at ({-\ndx},{\ndy}); \fi
  \ifnum\i=2 \coordinate (nb_\i) at ({+\ndx},{\ndy}); \fi
  \ifnum\i=3 \coordinate (nb_\i) at ({-\ndx},{0});    \fi
  \ifnum\i=4 \coordinate (nb_\i) at ({+\ndx},{0});    \fi
  \ifnum\i=5 \coordinate (nb_\i) at ({-\ndx},{-\ndy});\fi
  \ifnum\i=6 \coordinate (nb_\i) at ({+\ndx},{-\ndy});\fi
}
\foreach \i in {1,...,6}{
  \foreach \j in {1,...,6}{
    \ifnum\i<\j \draw[arr] (nb_\i) -- (nb_\j); \fi
  }
}
\foreach \k/\lbl in {1/3481\\群創,2/2301\\光寶科,3/2406\\國碩,4/3454\\晶豪,5/2390\\云辰,6/9912\\偉聯}{
  \node[circle, draw=thesisForest!80!thesisDark, fill=thesisForest!30, thick, minimum size=10mm]
       (Nb_\k) at (nb_\k) {};
  \draw[arr, looseness=7, out=105, in=75] (Nb_\k) to (Nb_\k);
  \node[nodelabel, below=5pt] at (Nb_\k.south) {\lbl};
}
\node[clustertitle] at (0,{1.35+0.75}) {光電及I/O};
\node[centerlabel] at (0, 2.4) {產業關係圖示（同產業完全連結與自迴圖結構）};
\end{scope}

%%── Cluster 3: 半導體 ─────────────────────────────────────────
\begin{scope}[xshift=10.4cm]
\foreach \i in {1,...,6}{
  \ifnum\i=1 \coordinate (nc_\i) at ({-\ndx},{\ndy}); \fi
  \ifnum\i=2 \coordinate (nc_\i) at ({+\ndx},{\ndy}); \fi
  \ifnum\i=3 \coordinate (nc_\i) at ({-\ndx},{0});    \fi
  \ifnum\i=4 \coordinate (nc_\i) at ({+\ndx},{0});    \fi
  \ifnum\i=5 \coordinate (nc_\i) at ({-\ndx},{-\ndy});\fi
  \ifnum\i=6 \coordinate (nc_\i) at ({+\ndx},{-\ndy});\fi
}
\foreach \i in {1,...,6}{
  \foreach \j in {1,...,6}{
    \ifnum\i<\j \draw[arr] (nc_\i) -- (nc_\j); \fi
  }
}
\foreach \k/\lbl in {1/2451\\創見,2/6799\\采頡,3/1560\\中砂,4/2337\\旺宏,5/3094\\聯傑,6/6271\\同欣電}{
  \node[circle, draw=thesisSand!80!thesisDark, fill=thesisSand!50, thick, minimum size=10mm]
       (Nc_\k) at (nc_\k) {};
  \draw[arr, looseness=7, out=105, in=75] (Nc_\k) to (Nc_\k);
  \node[nodelabel, below=5pt] at (Nc_\k.south) {\lbl};
}
\node[clustertitle] at (0,{1.35+0.75}) {半導體};
\end{scope}

\end{tikzpicture}
\end{document}
